% CREATED BY MAGNUS GUSTAVER, 2020
\chapter{Introduction}
\label{chap:introduction}

Large language model (LLM) agents are increasingly expected to operate across
long-running interactions rather than isolated prompts. A user may discuss a
task today, revise it next week, add a preference in a different session, and
later ask a question whose answer depends on several earlier moments. In this
setting, the current context window is not a sufficient representation of the
agent's state. The agent needs an external memory system that can preserve
relevant parts of past interactions and make them usable when a future query
arrives. Recent memory systems for LLM agents therefore study how to store,
organize, retrieve, and update long-term conversational information
\cite{chhikara2025mem0,li2025memos,xu2025amem,zhang2025gmemory,zhang2025memsim}.

This thesis studies a specific responsibility within such systems: the memory
module as an evidence provider. The memory module should not simply return a
long list of raw sessions, because that forces the downstream answer model to
perform retrieval, denoising, temporal reasoning, evidence selection, and final
answer generation at once. It should also not directly replace the answer model,
because then retrieval failures, evidence synthesis failures, and answer
generation failures become difficult to separate. Instead, a useful memory
module should return compact, faithful, source-grounded evidence that an answer
model can inspect and use. NanoMem is designed around this boundary. Its output
is structured evidence about what happened, when it happened, where it was
mentioned, and whether the accumulated evidence is sufficient for answering.

\begin{figure}[H]
  \centering
  \fbox{\begin{minipage}{0.9\textwidth}
    \centering
    \vspace{0.8em}
    \textcolor{red}{TODO: Replace this placeholder with the planned boundary
    diagram showing multi-session dialogue history, the memory evidence
    provider, structured evidence output, and the downstream answer model.}
    \vspace{0.8em}
  \end{minipage}}
  \caption{Role boundary for long-term agent memory. The memory system is
  treated as an evidence provider: it transforms dialogue history into compact
  structured evidence, while final answer generation remains outside the memory
  module.}
  \label{fig:intro-boundary}
\end{figure}

\section{Long-Term Conversational Memory}
\label{sec:intro-background}

Long-term conversational memory differs from ordinary document retrieval in two
ways. First, the source material is personal, incremental, and time-sensitive.
Sessions can contain preferences, plans, corrections, stale information,
future commitments, and references to earlier events. Second, the user often
does not remember the exact wording or location of the supporting evidence. A
query may refer to "the week after my promotion", "the restaurant I liked last
time", or "the book I was reading when the project changed direction". The
memory system must recover the relevant evidence even when the query does not
name the evidence directly.

Temporal reasoning is therefore central to agent memory. Benchmarks and prior
work show that LLMs remain brittle when they must reason over event order,
relative dates, and long conversational histories
\cite{chu2024timebench,fatemi2025testoftime,maharana2024locomo}. In a memory
setting, this difficulty is not only an answer-generation problem. It affects
which sessions are retrieved, which utterances are considered evidence, and
whether the system knows that more retrieval is needed. A memory system that
retrieves by semantic similarity alone may find a session mentioning a topic but
miss the actual event time. Conversely, a system that filters rigidly by session
timestamp may discard the session that contains the necessary temporal clue.

The example used throughout this thesis is a user asking: "What book was I
reading during the week I got promoted?" The query does not directly identify
the book session. A useful memory system must first locate evidence about the
promotion, infer the relevant week, and then use that inferred time window to
retrieve reading-related evidence. The first retrieval result is not the final
answer; it creates the search condition for the next retrieval step.

\section{The Evidence Addressability Gap}
\label{sec:intro-gap}

The central problem addressed in this thesis is the evidence addressability gap:
the evidence needed to answer a query is often not directly addressable from the
query text. The missing address may be a time window, an entity, a causal
predecessor, a remembered attribute, or a dependency exposed only after an
earlier retrieval step. This gap is broader than date normalization. It is a
problem of making hidden evidence searchable and inspectable.

Three coupled challenges make the gap difficult for agent memory systems.

The first challenge is the mismatch between session time and event time. A
session timestamp records when the user mentioned something. The event time
records when the event actually happened. If a user says on April 2 that they
"traveled to Europe last week", then April 2 is the mention time, while the trip
occurred during the previous week. A write-time memory system may extract the
travel event and attach the inferred event time, but doing only that loses the
fact that the user mentioned it on April 2. Storing both interpretations as
fixed memories can also confuse retrieval. The interpretation that matters is
often query dependent.

The second challenge is indirectly addressable temporal evidence. In the book
example, the book session becomes retrievable only after the promotion event has
been found and temporally grounded. A one-shot retriever receives the original
query but not the intermediate event-time inference. Even a strong reranker can
fail if the candidate set never contains the later evidence. The memory system
therefore needs a search state that can carry temporal hints forward across
retrieval rounds.

The third challenge is faithfulness under compression. Memory systems are often
used to reduce long histories into short evidence snippets. If a model compresses
uncertain cross-session reasoning into a statement that looks directly observed,
the downstream answer model cannot tell which part is supported by source text
and which part was inferred. For this reason, NanoMem separates source
identifiers, session time, event time, and inferred evidence in its structured
events. The goal is not only to find evidence, but also to preserve how the
evidence is grounded.

\begin{figure}[H]
  \centering
  \fbox{\begin{minipage}{0.9\textwidth}
    \centering
    \vspace{0.8em}
    \textcolor{red}{TODO: Replace this placeholder with the planned
    session-time/event-time and indirect-search diagram. The diagram should show
    why session-time search and single-pass temporal search miss different
    parts of the evidence trajectory.}
    \vspace{0.8em}
  \end{minipage}}
  \caption{Evidence addressability gap in temporal conversational memory. A
  session can mention an event at one time while the event happened at another
  time; after that event time is inferred, it may become the retrieval cue for a
  second piece of evidence.}
  \label{fig:intro-gap}
\end{figure}

\section{Limitations of Existing Memory Paradigms}
\label{sec:intro-existing-methods}

Existing memory systems approach long-term memory from several directions, each
of which addresses part of the problem but leaves the evidence addressability
gap unresolved.

Write-time memory methods extract, summarize, link, or update memory records
before the future query is known. Systems such as Mem0, MemOS, Zep, and A-Mem
make persistent memory more manageable by organizing user history into reusable
units or structured stores
\cite{chhikara2025mem0,li2025memos,rasmussen2025zep,xu2025amem}. This is useful
for storage and retrieval efficiency, but it forces the system to decide in
advance which interpretation of an utterance should be preserved. Future queries
may require a different temporal interpretation or a hidden dependency that was
not salient at ingestion time.

Retrieval and reranking methods keep more of the original history available at
query time. This avoids some premature compression, but it often returns raw
sessions or ranked snippets rather than structured evidence. The downstream
answer model must still determine which retrieved pieces are sufficient, which
temporal hints should be propagated, and whether another retrieval step is
needed. Temporal memory systems such as TReMu and Memory-T1 explicitly target
time-aware memory or temporal session selection
\cite{ge2025tremu,du2026memoryt1}, but they do not provide NanoMem's explicit
query-time evidence state with separate session-time and event-time fields.

Iterative or reflective retrieval systems are closer to the search behavior
needed here. Self-reflective retrieval and search-trained RAG systems can decide
when to retrieve again or adjust a query based on previous context
\cite{asai2024selfrag,jin2025searchr1,du2025memr3}. However, their intermediate
state is usually answer-oriented reasoning, retrieved snippets, or textual gap
descriptions. NanoMem instead makes the intermediate state an evidence pool:
structured event records with source identifiers, temporal grounding, and a
verdict about sufficiency.

\begin{table}[H]
  \centering
  \caption{Simplified positioning of NanoMem relative to common long-term memory
  paradigms. The detailed literature review is developed in Chapter 2.}
  \label{tab:intro-methods}
  \footnotesize
  \begin{tabular}{p{0.21\textwidth}p{0.23\textwidth}p{0.21\textwidth}p{0.23\textwidth}}
    \hline\hline
    Paradigm & Typical output & Strength & Limitation for this thesis \\
    \hline
    Write-time memory & Extracted records or graphs & Compact persistent state & Query-agnostic compression can miss future hidden dependencies \\
    Retrieval or reranking & Raw sessions, chunks, or ranked candidates & Preserves source text at query time & Does not synthesize sufficient structured evidence \\
    Iterative retrieval agent & Retrieval traces or answer reflections & Adapts search over multiple rounds & Loop state is rarely explicit temporal evidence \\
    NanoMem & Structured events and sufficiency verdict & Makes indirect temporal evidence addressable & Requires reliable evidence synthesis policy \\
    \hline\hline
  \end{tabular}
\end{table}

\section{Research Aim and Questions}
\label{sec:intro-rqs}

The aim of this thesis is to investigate whether a long-term conversational
memory system can construct compact, faithful, temporally grounded evidence at
query time, while exposing whether the evidence is sufficient for a downstream
answer model. This aim leads to three research questions.

\textbf{RQ1.} How can long-term conversational memory be represented as
source-grounded structured evidence, with explicit session time, event time, and
inference markers, rather than only as raw retrieved sessions or final answers?

\textbf{RQ2.} How can an iterative Planner-Synthesizer loop use evidence
constructed in earlier retrieval rounds to generate new retrieval cues and
recover temporal or causal evidence that is not named in the original query?

\textbf{RQ3.} How does NanoMem perform on long-term memory and temporal-causal
evidence-search benchmarks compared with write-time memory systems, retrieval
agents, temporal memory baselines, and full-context answer models, in terms of
answer quality, evidence compactness, temporal grounding, and sufficiency
judgment?

These questions define the structure of the thesis. RQ1 is addressed mainly by
the evidence schema and Temporal Evidence Pool in Chapter 3. RQ2 is addressed
by the iterative Planner-Retriever-Synthesizer loop and the sufficiency verdict.
RQ3 is addressed by the experimental evaluation in Chapter 4, including public
long-term memory benchmarks and ChainMem, a
diagnostic benchmark for indirectly addressable temporal-causal evidence.

\section{NanoMem Overview}
\label{sec:intro-nanomem}

NanoMem is a memory evidence synthesis framework for long-term conversational
agents. It combines lightweight temporal preprocessing, iterative retrieval, and
a trainable Synthesizer that produces structured evidence. At ingestion time,
relative temporal expressions in sessions are normalized into explicit temporal
hints where possible. This does not decide every future event interpretation in
advance; instead, it gives the query-time system better material for search.

At query time, a Planner proposes retrieval cues from the user query and the
current evidence state. A Retriever returns candidate sessions or chunks. A
Synthesizer then converts the retrieved material into evidence events and
updates the Temporal Evidence Pool. Each event records a source session, a
session time, an event time, and a concise evidence statement. The Synthesizer
also emits a verdict. A sufficient verdict stops retrieval and passes the
evidence pool to the downstream answer model. An insufficient verdict indicates
that the pool is relevant but incomplete, so the Planner should use the current
evidence and missing information to search again.

The Synthesizer is trained because prompting alone is not reliable enough for
the stop-or-continue decision. NanoMem uses supervised fine-tuning followed by
Group Relative Policy Optimization (GRPO) \cite{shao2024deepseekmath}. The
reward design combines parseable output format, compactness, downstream answer
usefulness, temporal grounding, and verdict accuracy. The verdict reward is
especially important because it supervises the internal search decision: when
the current evidence is incomplete, the correct behavior is not to answer
prematurely, but to preserve the partial evidence and continue retrieval.

\begin{figure}[H]
  \centering
  \fbox{\begin{minipage}{0.9\textwidth}
    \centering
    \vspace{0.8em}
    \textcolor{red}{TODO: Replace this placeholder with the planned NanoMem
    workflow diagram showing temporal preprocessing, Planner, Retriever,
    Synthesizer, Temporal Evidence Pool, verdict loop, and downstream answer
    model.}
    \vspace{0.8em}
  \end{minipage}}
  \caption{High-level NanoMem workflow. The Planner-Retriever-Synthesizer loop
  builds a Temporal Evidence Pool iteratively; the sufficiency verdict decides
  whether to continue retrieval or provide evidence to the answer model.}
  \label{fig:intro-workflow}
\end{figure}

\section{Contributions}
\label{sec:intro-contributions}

This thesis makes four contributions.

First, it formulates long-term agent memory as temporal-causal evidence
synthesis. The formulation separates the role of the memory system from the role
of the downstream answer model and focuses on evidence that is compact,
inspectable, and grounded in source sessions.

Second, it presents NanoMem, a Planner-Synthesizer memory framework that builds
a Temporal Evidence Pool through iterative retrieval. The evidence pool
explicitly separates event time from session time and preserves source
identifiers, allowing intermediate temporal inferences to guide subsequent
search.

Third, it trains the Synthesizer as an evidence construction policy. The
training procedure combines supervised warm-up and GRPO rewards for parseable
format, compactness, answer usefulness, temporal grounding, and sufficiency
judgment.

Fourth, it evaluates NanoMem on LoCoMo, LongMemEval, and ChainMem
\cite{maharana2024locomo,wu2024longmemeval}. The evaluation compares NanoMem
against write-time memory systems, search-time retrieval agents, temporal memory
baselines, and full-context answer models, with additional ablations and case
studies to analyze the role of temporal normalization, the evidence pool, and
the verdict signal.

\section{Scope and Delimitations}
\label{sec:intro-scope}

This thesis focuses on memory evidence construction, not on general final answer
generation. The downstream answer model is treated as an evidence consumer. This
boundary is intentional: it makes the memory module easier to inspect and makes
it possible to evaluate whether failures come from missing evidence, unfaithful
evidence synthesis, or final answer generation.

The temporal reasoning studied here is also scoped. NanoMem targets
multi-session conversational memory where relative time, event-time grounding,
and evidence discovered through intermediate temporal hints affect retrieval. It
does not claim to solve all forms of temporal logic, planning, or calendar
reasoning. The implemented system uses temporal normalization as a lightweight
preprocessing step, iterative retrieval over memory sessions or chunks, a
Synthesizer that produces structured evidence, and a separate downstream answer
model that consumes that evidence. The empirical claims in this thesis are
therefore limited to the LoCoMo, LongMemEval, and ChainMem evaluations and to
the concrete retriever settings, model backbones, answer models, and evaluation
protocols reported in Chapter 4.

\section{Thesis Structure}
\label{sec:intro-structure}

Chapter 2 reviews long-term agent memory, temporal reasoning in memory systems,
search-side retrieval, evidence synthesis, and reinforcement learning for
retrieval and memory agents.

Chapter 3 presents the NanoMem method. It defines the problem, the structured
evidence schema, the Temporal Evidence Pool, the Planner-Retriever-Synthesizer
loop, the verdict mechanism, the training objective, and the construction of
ChainMem.

Chapter 4 reports the experimental setup and results. It describes datasets,
baselines, metrics, main results, ablations, and qualitative case studies.

Chapter 5 discusses the implications and limitations of NanoMem, including
scalability, privacy, failure modes, and future work, before summarizing the
thesis.
